\cotefile{136}
\scanpage{235}
On considère l'anneau simplicial
\[
\widehat A_*^*(S,\gamma^*,t)^{\unsure{\wedge}}=\widehat A_*^{*\wedge}
\]
\unsure{par}
\[
\begin{aligned}
\widehat A_n^*&\simeq S\{\{X_0,\ldots,X_n\}\}/(\text{idéal }\inserted{\text{à p.d.}}\deleted{\text{\unsure{engendré}}}\inserted{\text{engendré par }\sum X_i-\unsure{t|\bar t}})\otimes_S\frac{\bigwedge^*[dX_0,\ldots,dX_n]}{\text{idéal }(\sum dX_i)}\\
(&\simeq S\{\{Y_0,\ldots,Y_{n-1}\}\}[dY_0,\ldots,dY_{n-1}]
\end{aligned}
\]
\deleted{$\simeq$}

On a donc
\[
\widehat A_*^*\simeq\widehat A_*^0\otimes_S\bigwedge^*\Psi_*^S.
\]
d'où
\[
\pi_\ell(\widehat A_*^i)=\pi_\ell(\widehat A_*^0\otimes\bigwedge^i\Psi_*^S)\simeq\pi_{\ell-i}(\widehat A_*^0)
\]
et il faut donc étudier $\pi_*(\widehat A_*^0)$.

On considère sur $\widehat A_*^0$ \deleted{$\unsure{\Gamma_S(\Phi_*)}$} la filtration décroissante par idéaux fermés stables par p.d. ``engendrée'' par la condition que \deleted{\unsure{pour}} \unsure{pour tout $n$}, les $X_i$ et $t$ sont en filtration $\geq1$. \unsure{On}
\[
\begin{aligned}
\mathrm{Fil}^0(\widehat A_n^0)&\simeq\deleted{\unsure{S}}\widehat A_n^0\\
\mathrm{Fil}^p(\widehat A_n^0)&\simeq\left\{\sum\deleted{\unsure{a}}a_{i_0\ldots i_n}\unsure{\big\{}X_0^{(i_0)}\cdots X_n^{(i_n)}\ \middle|\ a_{i_0\ldots i_n}\in S^{[p-\sum_0^n i_\alpha]}\right\}\\
&=\left\{\sum a_{i_0,\ldots i_{n-1}}Y_0^{(i_0)}\cdots Y_{n-1}^{(i_{n-1})}\ \middle|\ a_{i_0\ldots i_{n-1}}\in S^{\unsure{[p-\sum_0^{n-1}i_\alpha]}}\right\}.
\end{aligned}
\]
Bien sûr, la topologie \unsure{pr. div.} sur $\widehat A_n^0$ est aussi celle définie par cette filtration, et par suite $\widehat A_*^0$ est séparé et complet pour cette filtration. Un calcul immédiat donne, posant
\[
k=S/S^{[1]}
\]
\[
\mathrm{Gr}^p(\widehat A_*^0)\simeq\Gamma_k^p(\Phi_*^k)
\inserted{\simeq\widehat A_*^{p0}(\mathrm{Gr}^*(S),\widetilde\gamma^*,\widetilde t)}
\]
\margin{$\mathrm{Gr}^*(S)$ : $k$-\unsure{$\delta$}-algèbre graduée quotient de $k\unsure{\langle T\rangle|\{T\}}$.}
\inserted{si l'épi $k\unsure{\langle T\rangle|\{T\}}\longrightarrow\mathrm{Gr}^*(S)$ est un \uline{iso} \deleted{\unsure{et}} (donc $\mathrm{Gr}^p(\widehat A_*^0)$ flasque si $p\ne0$)}
d'où, dans cette hyp. \inserted{d'isomorphisme}
\[
\pi_i(\mathrm{Gr}^p(\widehat A_*^0))=\begin{cases}0&\text{si }p\ne0,\ i\ne0\ \deleted{\unsure{\longleftarrow}}\\k&\text{si }(i,p)=0\end{cases}
\]
d'où on conclut que \deleted{l'inclusion}
\[
\deleted{\mathrm{Fil}^1\widehat A_*^0\unsure{\hookrightarrow}\widehat A_*^0}
\]
\margin{234}
